\section{Discussion}\label{sec:discussion}

\subsection{A two-sided transfer result}

The combined result is sharper than either track alone.  Option D shows that
the M1.5 witness is not confined to the exact $(2,4)$ instance: it survives the
first alternative expansion while preserving the strongest structural control
available in this study.  Option C shows why that result should not be
generalized indiscriminately: increasing the voter dimension changes the
representation problem and introduces clauses that have no bundled
counterpart.

The resulting claim is directional:
\begin{itemize}
  \item the repair-presentation witness transfers cleanly along the
  two-voter alternative-expansion direction; and
  \item it does not automatically transfer along the voter-dimension direction.
\end{itemize}
This is not a dichotomy between success and failure.  The positive result
identifies a stable region, while the boundary result identifies the condition
that ends the \CM-protected comparison.

\subsection{What clause-multiset equivalence buys}

Clause-multiset equivalence is intentionally stronger than status agreement.
It creates a controlled comparison in which the same clauses are assigned to
different lever interfaces.  When the raw repair families differ under that
control, the difference is evidence about presentation granularity.

The control does not make one repair vocabulary canonical.  It shows the
opposite: the clause set alone does not select a unique human-facing repair
explanation.  A downstream report must decide whether to expose
person-specific levers, aggregate them under \texttt{minlib}, or present both
levels.  That decision is an abstraction contract and should be versioned and
audited like other parts of the pipeline.

\subsection{Limitations}

First, the evidence covers only $(2,5)$ and $(3,4)$ beyond the base case.  It is
not a statement for arbitrary electorate or alternative sizes.
Second, the rationality scope is local:
\texttt{no\_cycle3}$\wedge$\texttt{no\_cycle4} is not full
\texttt{SocialAcyclic}.  In particular, the $(2,5)$ result permits cycle lengths
outside the encoded short-cycle families.

Third, the Option D repair result is raw and lever-level.  It does not establish
that every semantic repair abstraction must differ, only that the current
bundled and split presentations differ before an explicit grouping policy is
applied.  Fourth, Option C repair enumeration is absent by design.  The
\texttt{sat\_equiv\_only} result leaves selector-strength explanations open, so
the paper records the boundary without speculating about an unmeasured repair
family.

Finally, the family comparisons are exploratory artifacts, not additions to the
baseline proof-carrying certificate set.  The Sen24 CNF remains a base-case
certificate and must not be cited as a family-level CNF proof.

\subsection{Implications for possibility atlases}

An atlas that reports only SAT/UNSAT status can treat the two Option C packages
as equivalent for a coarse boundary map.  An atlas that reports repairs cannot
stop there.  It needs metadata about the representation under which each repair
was computed and about the relation between compared encodings.

A practical report should therefore include at least:
\begin{enumerate}
  \item the active lever vocabulary;
  \item the repair-transport or grouping map;
  \item SAT/UNSAT status;
  \item the strongest verified structural relation, such as \CM;
  \item whether repair enumeration was authorized under that relation; and
  \item the scope caveat for the encoded rationality family.
\end{enumerate}
These fields prevent a status-equivalent boundary comparison from being
silently upgraded into a representation-invariant repair claim.

\subsection{Future decisions}

The next research decision is not implicit in this scaffold.  One path is to
freeze the two-sided M2.1 result and use it as the final claim boundary.
Another is to authorize a separate Option C repair comparison explicitly under
the weaker \texttt{sat\_equiv\_only} framing.  Such work would belong in an
exploratory appendix and would require wording that does not inherit the
Option D control.

Other extensions could test additional $(2,m)$ cases or design a bundled
encoding whose selector structure matches the Option C split closely enough
for a new \CM{} question.  Either extension would be a new experiment.  It is
not part of the evidence reported here.

\subsection{Conclusion}

Repair explanation is not a free corollary of impossibility transfer.  In the
Sen atlas, the base-case non-canonicity witness survives from $(2,4)$ to $(2,5)$
when the two-voter clause basis is preserved.  At $(3,4)$, UNSAT status still
survives, but the pair-selection and role machinery crosses a structural
boundary and breaks clause-multiset equivalence.  Treating these outcomes
together yields a disciplined conclusion: transfer of status, transfer of
clause structure, and transfer of repair presentation are separate engineering
claims and should be checked separately.
